%% ==============================================================================
%% figures/tikz/apercu-studio.tex
%% Illustration vectorielle d'un canevas de présentation moderne (standalone)
%% ==============================================================================
\documentclass[tikz, border=1mm]{standalone}

\usepackage{xcolor}
\providecolor{ThesisPrimary}{RGB}{15, 60, 110}
\providecolor{ThesisSecondary}{RGB}{20, 110, 120}
\providecolor{ThesisAccent}{RGB}{190, 75, 40}
\providecolor{ThesisHighlight}{RGB}{214, 134, 42}
\providecolor{ThesisText}{RGB}{30, 35, 42}
\providecolor{ThesisMuted}{RGB}{95, 105, 118}
\providecolor{ThesisBackground}{RGB}{246, 248, 251}

\usetikzlibrary{arrows.meta, shapes.geometric}

\begin{document}
\begin{tikzpicture}[font=\sffamily]
  % Cadre studio sombre façon interface Keynote / Figma
  \fill[ThesisPrimary!92!black, rounded corners=7pt] (0,0) rectangle (6.2, 4.1);
  % Barre supérieure minimaliste
  \fill[ThesisPrimary!75!black, rounded corners=7pt] (0,3.55) rectangle (6.2, 4.1);
  \fill[ThesisPrimary!75!black] (0,3.55) rectangle (6.2, 3.8);
  \fill[ThesisAccent] (0.35, 3.82) circle (2.6pt);
  \fill[ThesisHighlight] (0.58, 3.82) circle (2.6pt);
  \fill[ThesisSecondary] (0.81, 3.82) circle (2.6pt);
  \node[anchor=east, font=\tiny\bfseries\sffamily\color{white!75}] at (5.92, 3.82) {16:9 \textbullet{} VECTORIEL};

  % Vignette gauche (structure / rail)
  \fill[white!10, rounded corners=4pt] (0.28, 0.35) rectangle (1.65, 3.30);
  \foreach \y/\col in {2.85/ThesisAccent, 2.15/ThesisSecondary, 1.45/white!55, 0.75/white!35} {
    \fill[\col, rounded corners=2pt] (0.44, \y-0.16) rectangle (1.49, \y+0.26);
  }

  % Canevas principal simulé (carte claire haute résolution)
  \fill[white, rounded corners=4pt] (1.88, 0.35) rectangle (5.92, 3.30);
  \fill[ThesisAccent] (2.15, 2.95) rectangle (2.95, 3.04);
  \fill[ThesisText!85, rounded corners=1pt] (2.15, 2.66) rectangle (4.65, 2.82);
  \fill[ThesisMuted!50, rounded corners=1pt] (2.15, 2.44) rectangle (3.95, 2.54);

  % Mini graphique barres & courbe dans le canevas
  \fill[ThesisBackground, rounded corners=3pt] (2.15, 0.62) rectangle (4.05, 2.22);
  \foreach \x/\h/\c in {2.35/0.65/ThesisPrimary!45, 2.68/0.95/ThesisPrimary!70, 3.01/1.22/ThesisPrimary, 3.34/0.82/ThesisSecondary, 3.67/1.38/ThesisAccent} {
    \fill[\c, rounded corners=1.5pt] (\x, 0.78) rectangle (\x+0.22, 0.78+\h);
  }
  \draw[ThesisAccent, line width=1.4pt, smooth]
    plot coordinates {(2.46, 1.48) (2.79, 1.72) (3.12, 1.98) (3.45, 1.66) (3.78, 2.14)};
  \foreach \px/\py in {2.46/1.48, 3.12/1.98, 3.78/2.14} {
    \fill[white] (\px, \py) circle (2.2pt);
    \draw[ThesisAccent, line width=1pt] (\px, \py) circle (2.2pt);
  }

  % Colonne KPI droite dans la miniature
  \foreach \y/\col/\w in {1.88/ThesisPrimary/1.05, 1.28/ThesisSecondary/1.25, 0.68/ThesisAccent/0.92} {
    \fill[ThesisBackground, rounded corners=3pt] (4.22, \y-0.06) rectangle (5.68, \y+0.36);
    \fill[\col] (4.34, \y+0.15) circle (3.5pt);
    \fill[ThesisText!80, rounded corners=1pt] (4.52, \y+0.16) rectangle (4.52+\w*0.82, \y+0.24);
    \fill[ThesisMuted!55, rounded corners=1pt] (4.52, \y+0.03) rectangle (5.42, \y+0.10);
  }
\end{tikzpicture}
\end{document}
